% GENERATED FILE. Edit canonical lessons, then run npm run generate:books.
% canonical-source-hash: fnv1a64:dc3cae851ae36006
% canonical-lessons: FR-C58-voisin, FR-C58-invite, FR-C58-bebe, FR-C58-petit-fils, FR-C58-petite-fille

\chapter{Neighbour, Guest, Baby, Grandson, Granddaughter}
\label{ch:fr-neighbour-guest-baby-grandson-granddaughter}
\hlchaptermodality{58}

\begin{chapteropening}
\textbf{By the end of this chapter:} \emph{I can name a neighbour, a guest, a baby, a grandson, a granddaughter.}

Complete the last lesson of chapter 58: i can name a neighbour, a guest, a baby, a grandson, a granddaughter.
\end{chapteropening}

\section[le voisin]{le voisin — a neighbour}

\label{lesson:FR-C58-voisin}



\begin{quote}
Before the new one: say the French for an aunt, then the French for a cousin.
\end{quote}

\subsection*{You'll want to know: le voisin}
\textbf{le voisin} — \textquotedblleft{}a neighbour\textquotedblright{}.

\begin{grammarlens}[title={What you've built}]
One more word for this chapter.
\end{grammarlens}

\subsection*{Your turn}
\begin{itemize}
\raggedright
  \item \emph{Say it:} \emph{le voisin}
  \item \emph{Say it:} \emph{le voisin}, once more
  \item \emph{Say it:} say \emph{le voisin}
  \item \emph{Recall:} say the French for an aunt, then the French for a cousin, then say \emph{le voisin} again
\end{itemize}

\begin{tcolorbox}[breakable,colback=teal!4,colframe=teal!35!black,title={Before you move on}]
What does le voisin mean? (\textquotedblleft{}A neighbour\textquotedblright{}.) Say it once more.
\end{tcolorbox}
\section[l'invité]{l'invité — a guest}

\label{lesson:FR-C58-invite}



\begin{quote}
Before the new one: say the French for a cousin, then the French for a neighbour.
\end{quote}

\subsection*{You'll want to know: l'invité}
\textbf{l'invité} — \textquotedblleft{}a guest\textquotedblright{}.

\begin{grammarlens}[title={What you've built}]
One more word for this chapter.
\end{grammarlens}

\subsection*{Your turn}
\begin{itemize}
\raggedright
  \item \emph{Say it:} \emph{l'invité}
  \item \emph{Say it:} \emph{l'invité}, once more
  \item \emph{Say it:} say \emph{l'invité}
  \item \emph{Recall:} say the French for a cousin, then the French for a neighbour, then say \emph{l'invité} again
\end{itemize}

\begin{tcolorbox}[breakable,colback=teal!4,colframe=teal!35!black,title={Before you move on}]
What does l'invité mean? (\textquotedblleft{}A guest\textquotedblright{}.) Say it once more.
\end{tcolorbox}
\section[le bébé]{le bébé — a baby}

\label{lesson:FR-C58-bebe}



\begin{quote}
Before the new one: say the French for a neighbour, then the French for a guest.
\end{quote}

\subsection*{You'll want to know: le bébé}
\textbf{le bébé} — \textquotedblleft{}a baby\textquotedblright{}.

\begin{grammarlens}[title={What you've built}]
One more word for this chapter.
\end{grammarlens}

\subsection*{Your turn}
\begin{itemize}
\raggedright
  \item \emph{Say it:} \emph{le bébé}
  \item \emph{Say it:} \emph{le bébé}, once more
  \item \emph{Say it:} say \emph{le bébé}
  \item \emph{Recall:} say the French for a neighbour, then the French for a guest, then say \emph{le bébé} again
\end{itemize}

\begin{tcolorbox}[breakable,colback=teal!4,colframe=teal!35!black,title={Before you move on}]
What does le bébé mean? (\textquotedblleft{}A baby\textquotedblright{}.) Say it once more.
\end{tcolorbox}
\section[le petit-fils]{le petit-fils — a grandson}

\label{lesson:FR-C58-petit-fils}



\begin{quote}
Before the new one: say the French for a guest, then the French for a baby.
\end{quote}

\subsection*{You'll want to know: le petit-fils}
\textbf{le petit-fils} — \textquotedblleft{}a grandson\textquotedblright{}.

\begin{grammarlens}[title={What you've built}]
One more word for this chapter.
\end{grammarlens}

\subsection*{Your turn}
\begin{itemize}
\raggedright
  \item \emph{Say it:} \emph{le petit-fils}
  \item \emph{Say it:} \emph{le petit-fils}, once more
  \item \emph{Say it:} say \emph{le petit-fils}
  \item \emph{Recall:} say the French for a guest, then the French for a baby, then say \emph{le petit-fils} again
\end{itemize}

\begin{tcolorbox}[breakable,colback=teal!4,colframe=teal!35!black,title={Before you move on}]
What does le petit-fils mean? (\textquotedblleft{}A grandson\textquotedblright{}.) Say it once more.
\end{tcolorbox}
\section[la petite-fille]{la petite-fille — a granddaughter}

\label{lesson:FR-C58-petite-fille}



\begin{quote}
Before the new one: say the French for a baby, then the French for a grandson.
\end{quote}

\subsection*{You'll want to know: la petite-fille}
\textbf{la petite-fille} — \textquotedblleft{}a granddaughter\textquotedblright{}.

\begin{grammarlens}[title={What you've built}]
One more word for this chapter.
\end{grammarlens}

\subsection*{Your turn}
\begin{itemize}
\raggedright
  \item \emph{Say it:} \emph{la petite-fille}
  \item \emph{Say it:} \emph{la petite-fille}, once more
  \item \emph{Say it:} say \emph{la petite-fille}
  \item \emph{Recall:} say the French for a baby, then the French for a grandson, then say \emph{la petite-fille} again
\end{itemize}

\begin{tcolorbox}[breakable,colback=teal!4,colframe=teal!35!black,title={Before you move on}]
What does la petite-fille mean? (\textquotedblleft{}A granddaughter\textquotedblright{}.) Say it once more.
\end{tcolorbox}
